% Figure: FBD of the upper wall anchor C of the jib crane, isolating the
% bracket plate so the bolt group load path is explicit.
\begin{figure}[htbp]
\centering
\begin{tikzpicture}[
  load/.style={draw=engred, -{Stealth}, very thick},
  tie/.style={draw=engblue, -{Stealth}, very thick},
  scale=1.0,
]
% wall
\draw[draw=engblack, very thick] (0,-0.4) -- (0,3.6);
\foreach \y in {-0.2,0.2,...,3.4} {\draw[draw=enggray, thin] (0,\y) -- (-0.26,\y-0.18);}
% anchor plate
\draw[draw=engblack, fill=engblue!10, thick] (0,0.6) rectangle (0.7,2.6);
\node[font=\scriptsize, anchor=west] at (0.75,2.3) {anchor plate};
% two bolts into the wall
\draw[draw=engblack, fill=engamber!50] (0.05,1.0) rectangle (0.35,1.25);
\draw[draw=engblack, fill=engamber!50] (0.05,1.95) rectangle (0.35,2.2);
\node[font=\scriptsize, anchor=west] at (0.4,1.12) {bolt 2};
\node[font=\scriptsize, anchor=west] at (0.4,2.07) {bolt 1};
% tie-rod pin at the plate (point C)
\coordinate (C) at (0.7,1.6);
\fill[engblack] (C) circle (0.06);
\node[font=\scriptsize, anchor=south west] at (C) {$C$};
% the tie rod pulls the plate down and to the right (reaction to T): the
% rod runs down to the boom tip, so on the anchor it pulls toward the tip
\draw[tie] (C) -- ++(2.0,-1.2) node[anchor=west, font=\small, engblue] {$T$};
% bolt tension on the top bolt (the prying couple) and bearing at bottom
\draw[load] (0.2,2.2) -- ++(-1.3,0) node[anchor=east, font=\scriptsize, engred] {$T_{b1}$};
\draw[load] (-1.0,1.0) -- ++(1.3,0) node[anchor=west, font=\scriptsize, engred] {$C_{b2}$};
\end{tikzpicture}
\caption[Free-body diagram of the jib-crane wall anchor]{The upper wall
anchor $C$ isolated as its own body. The tie rod pulls on the plate with
force $T$, directed toward the boom tip. The plate is held by a two-bolt
group: the moment of the off-axis pull about the bottom edge is resisted
as a couple, putting the top bolt in tension $T_{b1}$ and the bottom of
the plate in bearing $C_{b2}$ against the wall. Tracing the load through
the plate to the individual bolts is the step that decides whether the
anchor or the tie rod governs, and it is the step the original Hyatt
connection skipped.}
\label{fig:vol03ch01:jib-anchor}
\end{figure}
